\section{Swarm}
As mentioned earlier, cost constraints meant we chose a non-HA Swarm cluster setup.

The previous CI/CD setup redeployed the entire Docker Compose stack, including the database, monitoring services, and application, causing brief downtime during updates. It also lacked rollback resilience: if a new Docker image failed to start, the application could not recover to the previous working version, potentially causing extended downtime.
To solve this, we migrated the application to Docker Swarm. Swarm supports individual container scaling and deployment control, allowing the main application to run across two containers behind a built-in round-robin load balancer.

When a new image is deployed, the CI/CD pipeline notifies the Swarm cluster. Swarm then updates containers one at a time, waiting for each new container to become healthy before replacing the next. Other services, such as the database and monitoring stack, remain untouched during the process. This effectively eliminates downtime for end users during deployments.

% Add keywords to start and stop
